\chapter{Introducao}
\label{chap:n1_intro}

\textcolor{red}{Senti falta nos paragrafos seguintes de algumas referencias mais basicas como uma geral para avaliação de clustering, uma mais específica sobre indices internos e externos, outra para ARI, uma geral para TRI, uma ou mais para TRI aplicado a AM}

\textcolor{blue}{Correção: adicionadas referências gerais para avaliação de agrupamentos e índices internos e externos \cite{wu2020decision,WOS:000172252500002,WOS:000232704100004,deborah2010survey}, para o ARI \cite{hubert1985comparing}, para TRI \cite{embretson2013item} e para TRI aplicada a AM \cite{chen2019birt, martinez2016making, moraes2020item, moraes2022evaluating}. Pendente: uma referência de \textit{survey} mais específica sobre a distinção entre índices internos e externos.}

A avaliação de modelos é um dos pilares centrais do \gls{AM}, orientando a seleção de modelos, a interpretação de resultados e o avanço metodólogico da área. No contexto supervisionado, métricas de desempenho podem ser diretamente calculadas a partir de rótulos verdadeiros, uma vez que são conhecidos, permitindo comparações mais objetivas entre desempenho de modelos. Entretanto, no aprendizado não supervisionado, particularmente em tarefas de agrupamento, a ausência de rótulos reais torna o problema de avaliação significativamente mais desafiador. Nesses cenários, a qualidade de uma partição deve ser inferida a partir da própria estrutura dos dados ou da comparação indireta entre modelos, levantando questões sobre robustez, interpretabilidade e generalidade das medidas utilizadas.

Tradicionalmente, a validação de agrupamentos é realizada por meio de medidas internas ou externas \cite{wu2020decision,WOS:000172252500002,WOS:000232704100004,deborah2010survey}. Medidas externas como o \textit{Adjusted Rand Index} \cite{hubert1985comparing} ajustado dependem da existência de uma partição de referência e são amplamente utilizadas em experimentos controlados. Contudo, tais medidas tornam-se inviáveis em aplicações reais, onde o objetivo do agrupamento desconhece os rótulos verdadeiros. Medidas internas, por sua vez, avaliam propriedades como coesão e separação dos clusters, mas são sensíveis à métrica de distância adotada, podendo apresentar vieses em relação a determinadas estruturas e frequentemente produzem ranqueamentos contraditórios entre modelos. Além disso, ambas as abordagens negligenciam uma dimensão fundamental: a variabilidade de dificuldade entre instâncias, tratando todos os pontos de dados com uma dificuldade constante.

Nesse contexto, esta dissertação adota a \gls{TRI} \cite{embretson2013item} como arcabouço teórico para modelar avaliação no cenário de \gls{AM}. Originalmente desenvolvida na psicometria para estimar habilidades latentes de respondentes e dificuldades de itens, a \gls{TRI} permite analisiar desempenho considerando explicitamente a interação entre habilidade e dificuldade. Em anos recentes, modelos baseados em TRI foram adaptados para avaliar sistemas de Inteligência Artificial em tarefas supervisionadas \cite{chen2019birt, martinez2016making, moraes2020item, moraes2022evaluating}, possibilitando estimar habilidade de modelos e identificar instâncias mais desafiadoras ou podendo ser consideradas como ruidosas. Essa perspectiva oferece uma alternativa às métricas tradicionais, ao introduzir uma modelagem probabilística que incorpora um disciminante e uma dificuldade associadas as instâncias, fornecendo interpretações mais refinadas sobre o desempenho.

\textcolor{red}{@Ricardo:eu tiraria o paragrafo seguinte}

\textcolor{blue}{Correção: parágrafo sobre a contribuição do $\beta^{4}$-\gls{IRT} removido.}

Nessa dissertação, propomos o método entitulado \gls{CLAIRE}, que reformula a avaliação de agrupamentos como um problema de estimação de habilidade latente via \gls{TRI}. Em vez de depender do rótulo real da instância, o método constrói uma matriz de resposta baseada na concordância entre modelos quanto ao agrupamento de pares de instâncias. Sob essa perspectiva, modelos mais informativos são aqueles que tendem a concordar entre si nas instâncias mais difíceis, enquanto instâncias difíceis são aquelas nas quais apenas modelos de alta habilidade apresentam um consenso. Experimentos demonstram que essa abordagem produz ranqueamentos altamente correlacionados com medidas externas tradicionais e permanece robusta mesmo quando partições aleatórias são adicionadas ao conjunto de modelos avaliados.


Essa formulação introduz uma mudança conceitual relevante: uma avaliação deixa de ser uma função puramente baseado nos dados ou uma comparação direta com um rótulo externo, passando a ser uma propriedade baseada na interação entre modelos e instâncias. Como consequência, a qualidade de um modelo é definida não apenas pela consistência global de sua partição, mas pela sua capacidade de manter concordância em regiões mais desafiadoras no espaço do conjunto de dados. Essa perspectiva também possiblita análises adicionais, como identificação de regiões de incerteza, detecção de instãncias ambíguas e estudo de impacto de modelos fracos no processo de avaliação.

Nos experimentos realizados, o método \gls{CLAIRE} \cite{ferreira2025claire} foi utilizado para avaliar 35 diferentes métodos de agrupamento produzidos por 6 diferentes algoritmos de agrupamento conhecidos na literatura. O modelo $\beta^{4}$-\gls{IRT} \cite{ferreirajunior2023beta4irtnewbeta3irtenhanced} foi utilizado para estimar habilidades e dificuldades para cada conjunto de dados trabalhado. Os experimentos realizados foram sobre 9 conjuntos de dados com diferentes caracteristicas (forma, sobreposição de classes e instâncias ruidosas). Os resultados revelaram uma alta correlação entre habilidade e 6 diferentes métricas de validação, das quais 3 eram externas e 3 internas. Esse resultado indica que a habilidade foi suficiente para identificar modelos que são capazes de produzir boas partições dos dados. Essa contribuição do artigo pode ser resumida em alguns tópicos, sendo eles:

\begin{itemize}
    \item Uma nova abordagem para validação de métodos não supervisionados, em especial agrupamentos, baseado em \gls{TRI} e Concordância entre Agrupamentos, abrindo novas possibilidades para avaliação de algoritmos em outros cenários não supervisionados ou semi-supervisionados, aonde o rótulo real é desconhecido;

    \item O método \gls{CLAIRE} não depende do número de cluster a priori, e pode ser utilizado não apenas para validação de agrupamentos, mas também para estimação de dificuldade de instâncias, quantificação de regiões de incerteza em um conjunto de dados;
    
    \item Experimentos realizados revelaram uma alta correlação  da habilidade estimada e métricas de validação, também mostrando uma forte robustez à presença de modelos fracos, dentro do grupo utilizados, tornando-se bastante promissor,ma vez que a solução não necessita de informações externas para uma avaliação precisa.

\end{itemize}

Ao integrar avanços metodológicos na modelagem de discriminação com uma nova fomulação para validação de agrupamentos, esta dissertação estabelece a habilidade latente como princípio unificador para avaliação de métodos não supervisionados. Dessa forma, contribui tanto para o desenvolvimento teórico de métricas de avaliação quanto para aplicações práticas em cenários onde a ausência de rótulos é a regra. A abordagem proposta não se limita a introduzir novos métodos, mas estabelece uma estrutura probabilística para avaliação no cenário não supervisionado, na qual habilidade, dificuldade e discriminação são quantidades modeladas e interpretáveis. Essa estrutura amplia o entendimento sobre o comportamento de métodos de agrupamento, abrindo caminho para extensões futuras em outras tarefas não supervisionadas.


 
 